\chapter{Martingales}

There is a particular fantasy that has bankrupted more people than any other in the history of
gambling. It goes like this. You walk up to a roulette table and bet one dollar on red. If you win,
great --- pocket the dollar and start over. If you lose, you simply double your bet to two dollars.
Lose again? Bet four. Lose again? Bet eight. The first time red finally comes up, you recover
everything you have lost \emph{plus} one dollar of profit. Since red is essentially certain to appear
\emph{eventually}, you have apparently discovered a way to extract guaranteed money from a game with
no edge. This strategy is so famous it has a name: the \textbf{martingale}.

It does not work. And the reason it does not work --- the reason \emph{nothing} works against a fair
game --- is one of the deepest facts in probability. To understand it, we will spend this chapter
building the mathematical object that, by an irony worth savoring, inherited the name of the betting
system it destroys: the martingale.

A martingale is the precise mathematical formalization of a \emph{fair game}: a running fortune whose
expected next value, given everything you currently know, is exactly its present value. ``On average,
you neither gain nor lose.'' That one sentence, taken seriously, carries a long way. It lets
us compute gambler's-ruin probabilities in a single line, explains why no clever stopping rule can
beat a fair casino, prices financial options, and underpins a whole convergence theory that rivals the
law of large numbers in importance. This chapter leans hard on conditional expectation, so if
$\E[Y \given X]$ and the tower property feel shaky, glance back at that chapter before we
begin --- everything here is conditional expectation wearing a story.

\section{Fair games and the definition}
\label{ch28:definition}

Imagine watching a quantity evolve in time: your bankroll at a casino, the price of a stock, the
position of a wandering particle, your running estimate of some unknown parameter as data trickles in.
Call its value at time $n$ the random variable $M_n$, for $n = 0, 1, 2, \dots$. At each moment you
have observed a \emph{history} --- everything that has happened up through time $n$. We bundle that
history into the symbol $\mathcal{F}_n$, read ``the information available at time $n$.'' As time
passes you learn more, never less, so $\mathcal{F}_0 \subseteq \mathcal{F}_1 \subseteq \mathcal{F}_2
\subseteq \cdots$. This growing family of information is called a \textbf{filtration}, and you should
picture it as the expanding contents of your own memory.

\begin{intuitionbox}
Do not be intimidated by $\mathcal{F}_n$. In every concrete example in this chapter it means nothing
more exotic than ``the values of $M_0, M_1, \dots, M_n$ that you have already seen,'' possibly together
with some other observed variables. Conditioning on $\mathcal{F}_n$ just means ``given everything we
know so far.'' When you read $\E[M_{n+1} \given \mathcal{F}_n]$, hear it as: ``my best forecast of
tomorrow's fortune, using my entire current memory.''
\end{intuitionbox}

What does it mean for the game to be \emph{fair}? It means that knowing the entire past, my best guess
for tomorrow's fortune is exactly today's fortune. I expect to walk away with what I came with. That is
the whole definition.

\begin{definition}[Martingale]
\label{ch28:def-mtg}
A sequence of random variables $(M_n)_{n \ge 0}$ adapted to a filtration $(\mathcal{F}_n)_{n\ge 0}$
(meaning $M_n$ is determined by $\mathcal{F}_n$) with $\E\abs{M_n} < \infty$ for every $n$ is a
\emph{martingale} if
\[
  \E[M_{n+1} \given \mathcal{F}_n] \;=\; M_n \qquad \text{for all } n \ge 0.
\]
If instead $\E[M_{n+1} \given \mathcal{F}_n] \ge M_n$ for all $n$, the sequence is a
\emph{submartingale} (a favorable game, drifting up); if $\E[M_{n+1} \given \mathcal{F}_n] \le M_n$, it
is a \emph{supermartingale} (an unfavorable game, drifting down).
\end{definition}

\begin{trapbox}
The ``super/sub'' naming feels backwards to everyone at first. A \emph{super}martingale goes
\emph{down} (it is the gambler's view of a casino: the house edge drags your fortune down). A
\emph{sub}martingale goes \emph{up}. The mnemonic that sticks: a supermartingale is ``superior for the
house,'' a submartingale is ``subpar for the house.'' The deeper reason is an analogy with harmonic
functions, where superharmonic functions lie above their spherical averages, so the expectation drops.
\end{trapbox}

The integrability condition $\E\abs{M_n} < \infty$ is not bureaucratic fine print; it is what makes
the conditional expectations exist as genuine finite numbers. We will keep it in the background but it
is always quietly assumed.

\subsection{The first and most important consequence}

Take the defining equation and average out the conditioning. By the tower property of conditional
expectation,
\[
  \E[M_{n+1}] \;=\; \E\big[\E[M_{n+1}\given\mathcal{F}_n]\big] \;=\; \E[M_n].
\]
So the expectation never moves: $\E[M_n] = \E[M_{n-1}] = \cdots = \E[M_0]$ for every $n$. A
martingale has \textbf{constant expectation}. This is the seed from which the entire chapter grows.

\begin{focusbox}
\textbf{The one fact to remember.} For a martingale, $\E[M_n] = \E[M_0]$ for all $n$. No matter how
long you play, your expected fortune equals your starting fortune. The drama of the rest of the chapter
is about whether you can cheat this by choosing a clever \emph{time to stop}. (Spoiler: under mild
conditions, you cannot.)
\end{focusbox}

\begin{historybox}
The word \emph{martingale} originally named a piece of horse tack --- a strap preventing a horse from
throwing its head back --- and by the 18th century it had attached itself to the doubling betting
system described in this chapter's opening. The mathematical object was named after the strategy by
Jean Ville in 1939, who used it to repair Richard von Mises's frequentist theory of randomness. But the
theory truly came alive in the hands of Joseph Leo Doob, whose 1953 book \emph{Stochastic Processes}
turned martingales into one of the central tools of modern probability. Doob's optional stopping
theorem and martingale convergence theorem are load-bearing pillars of everything from stochastic
calculus to the analysis of algorithms. It is a charming accident of history that the deepest theorem
about why you \emph{cannot} beat a fair game is named after the betting system that promised you could.
\end{historybox}

\section{A gallery of martingales}
\label{ch28:examples}

Definitions only become real when you have seen enough examples that the pattern becomes visible. Here
are the canonical ones; each will return later.

\begin{example}[Symmetric random walk --- the prototype]
\label{ch28:ex-walk}
Let $X_1, X_2, \dots$ be independent with $\P(X_i = +1) = \P(X_i = -1) = \tfrac12$, and set
$S_0 = 0$, $S_n = X_1 + \cdots + X_n$. Take $\mathcal{F}_n$ to be the history of the first $n$ steps.
Then
\[
  \E[S_{n+1}\given\mathcal{F}_n] = \E[S_n + X_{n+1}\given\mathcal{F}_n]
  = S_n + \E[X_{n+1}] = S_n + 0 = S_n.
\]
So $(S_n)$ is a martingale. This is the fair coin-flipping game: bet a dollar on each flip, win or lose
a dollar, and your fortune wanders with no drift. If the coin is biased toward losing
($\P(X_i = +1) = p < \tfrac12$), then $\E[S_{n+1}\given \mathcal{F}_n] = S_n + (2p-1) < S_n$ and you
have a supermartingale --- the accurate model of a real casino.
\end{example}

Figure~\ref{ch28:fig-walk} shows a few sample paths of this fair game. Each wanders without drift, and
although individual paths roam far from zero, the average across paths stays pinned at the starting
value --- the visible signature of constant expectation.

\begin{figure}[H]
\centering
\begin{tikzpicture}
\begin{axis}[
    width=12cm, height=6.0cm,
    xlabel={step $n$}, ylabel={fortune $S_n$},
    title style={font=\small},
    title={Symmetric random walk: a fair game wanders with no drift},
    grid=both, grid style={gray!20},
    xmin=0, xmax=40, ymin=-9, ymax=9,
    legend pos=outer north east, legend style={font=\footnotesize},
]
\addplot[thick, sourceblue] coordinates {
(0,0)(2,2)(4,0)(6,2)(8,4)(10,2)(12,4)(14,6)(16,4)(18,6)(20,4)(24,6)(28,4)(32,6)(36,8)(40,6)};
\addlegendentry{path A}
\addplot[thick, targetgreen] coordinates {
(0,0)(2,-2)(4,-4)(6,-2)(8,-4)(10,-2)(12,0)(14,-2)(16,-4)(18,-2)(20,-4)(24,-2)(28,0)(32,-2)(36,0)(40,-2)};
\addlegendentry{path B}
\addplot[thick, quotientorange] coordinates {
(0,0)(2,2)(4,0)(6,-2)(8,0)(10,2)(12,0)(14,2)(16,0)(18,-2)(20,0)(24,2)(28,0)(32,2)(36,0)(40,2)};
\addlegendentry{path C}
\addplot[densely dashed, trapred, thick] coordinates {(0,0)(40,0)};
\addlegendentry{$\E[S_n]=0$}
\end{axis}
\end{tikzpicture}
\caption{Three realizations of the symmetric simple random walk $S_n = X_1 + \cdots + X_n$. No path
shows any systematic tendency to rise or fall; the expected fortune (dashed) stays exactly at the
starting value of zero for all $n$. This is what ``fair'' looks like: motion without drift.}
\label{ch28:fig-walk}
\end{figure}

\begin{example}[The gambler's variance martingale]
\label{ch28:ex-variance}
With the same symmetric walk, consider $M_n = S_n^2 - n$. Watch what happens:
\[
  \E[S_{n+1}^2 \given \mathcal{F}_n] = \E[(S_n + X_{n+1})^2 \given \mathcal{F}_n]
  = S_n^2 + 2 S_n \E[X_{n+1}] + \E[X_{n+1}^2] = S_n^2 + 1,
\]
since $\E[X_{n+1}] = 0$ and $X_{n+1}^2 = 1$ always. Therefore
$\E[S_{n+1}^2 - (n+1)\given\mathcal{F}_n] = S_n^2 + 1 - (n+1) = S_n^2 - n = M_n$. So $S_n^2 - n$ is a
martingale even though $S_n^2$ itself drifts upward at a steady rate of one per step. The ``$-n$''
exactly cancels the predictable growth of the variance. Later we use this object to compute
\emph{how long} a gambler's-ruin game lasts.
\end{example}

\begin{example}[Likelihood ratios]
\label{ch28:ex-lr}
Here is a martingale that does not look like gambling at all but powers all of sequential statistics.
Suppose data $Y_1, Y_2, \dots$ are i.i.d.\ from some true density $f$, but you entertain a rival
hypothesis with density $g$. Form the running \emph{likelihood ratio}
\[
  L_n = \prod_{i=1}^{n} \frac{g(Y_i)}{f(Y_i)}, \qquad L_0 = 1.
\]
Then, because $Y_{n+1}$ is independent of the past and integrates against the \emph{true} density $f$,
\[
  \E[L_{n+1}\given\mathcal{F}_n] = L_n \cdot \E\!\left[\frac{g(Y_{n+1})}{f(Y_{n+1})}\right]
  = L_n \int \frac{g(y)}{f(y)} f(y)\,\dd y = L_n \int g(y)\,\dd y = L_n.
\]
So under the true model $f$, the likelihood ratio is a martingale with mean $1$ forever --- even though
it is built from products, not sums. This single fact is the engine behind Wald's sequential
probability ratio test, and the martingale convergence theorem will tell us that $L_n$ converges (in
fact, to $0$ almost surely when $f \ne g$, which is exactly why hypothesis testing eventually
discriminates between the two models).
\end{example}

\begin{example}[Polya's urn]
\label{ch28:ex-polya}
An urn starts with one red and one black ball. At each step you draw a ball uniformly at random,
return it, \emph{and add one more ball of the same color}. Rich get richer. Let $R_n$ be the number of
red balls after $n$ draws and let $M_n = R_n / (n+2)$ be the \emph{fraction} of red balls (there are
$n+2$ balls total after $n$ draws). Given the current composition, the chance the next ball is red is
exactly the current red fraction $M_n$, so
\[
  \E[R_{n+1}\given\mathcal{F}_n] = R_n + M_n = R_n + \frac{R_n}{n+2} = R_n \cdot \frac{n+3}{n+2},
\]
and dividing by $n+3$ gives $\E[M_{n+1}\given\mathcal{F}_n] = R_n/(n+2) = M_n$. The red fraction is a
martingale! This is genuinely surprising: a process with strong positive feedback nonetheless keeps its
proportion ``fair'' in expectation. By the convergence theorem the fraction settles to a random limit
--- and that limit turns out to be uniformly distributed on $[0,1]$, which is a surprising fact in its own right.
\end{example}

\begin{example}[Doob martingale: revealing information]
\label{ch28:ex-doob}
Let $Z$ be any integrable random variable and let information arrive in pieces $\mathcal{F}_0 \subseteq
\mathcal{F}_1 \subseteq \cdots$. Define $M_n = \E[Z \given \mathcal{F}_n]$, your best estimate of $Z$
given what you know at time $n$. The tower property instantly gives
$\E[M_{n+1}\given\mathcal{F}_n] = \E[\E[Z\given\mathcal{F}_{n+1}]\given\mathcal{F}_n]
= \E[Z\given\mathcal{F}_n] = M_n$. So \emph{any} sequence of sharpening estimates of a fixed quantity
is automatically a martingale. This ``Doob martingale'' is the backbone of concentration inequalities
(McDiarmid's inequality) used everywhere in machine learning to bound how far a complicated function of
many independent inputs can stray from its mean.
\end{example}

\begin{happeningbox}
Notice the unifying thread. A martingale is what you get whenever ``the best prediction of the next
value, given the past, is the current value.'' Sums of fair bets, products of mean-one factors, and
running conditional expectations of a fixed target all share this property. The deep reason is that the
operation $Z \mapsto \E[Z\given\mathcal{F}_n]$ is a \emph{projection}: it throws away the
unpredictable part and keeps only what the current information can see. Martingale theory is, at heart,
the dynamics of these projections as information grows. (This is also why the rigorous chapter on
conditional expectation later in the book treats it as an orthogonal projection in a space of random
variables --- the geometry is real.)
\end{happeningbox}

\section{Why you cannot beat a fair game}
\label{ch28:cannot-beat}

Now we confront the gambler's dream head-on. Surely \emph{some} betting strategy can turn a fair game
into a winning one? You get to choose how much to bet each round based on everything you have seen. Can
cleverness manufacture an edge?

Let us formalize ``a strategy.'' Before round $n$, knowing the history $\mathcal{F}_{n-1}$, you choose
a stake $C_n$ --- how many units to wager. The crucial constraint is that $C_n$ may depend only on the
\emph{past}, not on the future outcome of round $n$ (no clairvoyance, no peeking). A sequence with this
property, $C_n$ determined by $\mathcal{F}_{n-1}$, is called \textbf{predictable}. Your winnings on
round $n$ are $C_n(M_n - M_{n-1})$, so your fortune after $n$ rounds is the \emph{discrete stochastic
integral}, or \emph{martingale transform},
\[
  W_n = W_0 + \sum_{k=1}^{n} C_k\,(M_k - M_{k-1}).
\]

\begin{theorem}[You cannot beat a fair game]
\label{ch28:thm-transform}
If $(M_n)$ is a martingale and $(C_n)$ is a bounded predictable sequence, then the transformed process
$(W_n)$ above is also a martingale. Consequently $\E[W_n] = W_0$ for every $n$: no predictable betting
strategy changes your expected fortune.
\end{theorem}

\begin{whybox}
The proof is one line of conditional expectation, and it is worth savoring because it is the
mathematical core of ``the house always wins on average.'' Condition on the past $\mathcal{F}_{n-1}$.
The stake $C_n$ is already known --- it was chosen before the round --- so it slides out of the
conditional expectation as a known constant:
\[
  \E[W_n - W_{n-1}\given\mathcal{F}_{n-1}]
  = \E[C_n(M_n - M_{n-1})\given\mathcal{F}_{n-1}]
  = C_n\,\E[M_n - M_{n-1}\given\mathcal{F}_{n-1}]
  = C_n \cdot 0 = 0.
\]
The last step is the martingale property itself: $\E[M_n \given \mathcal{F}_{n-1}] = M_{n-1}$. So each
round contributes zero expected change \emph{no matter how you size your bet}, because you sized it
before you knew the outcome. The fairness of the underlying game leaks through any strategy built only
from the past. The only way to win would be to let $C_n$ depend on the result of round $n$ --- to bet
after seeing the outcome --- which is precisely cheating.
\end{whybox}

This already rules out infinitely many ``systems.'' Doubling, the Fibonacci system, anti-martingales,
trend-following --- all of them are predictable strategies, so all of them leave your expected fortune
untouched. But the doubling system seemed to promise a \emph{guaranteed} win, not merely a fair
expectation. To finish it off we need to think about \emph{stopping times}, and that brings us to the
crown jewel of the chapter.

\section{Stopping times and optional stopping}
\label{ch28:optional-stopping}

A gambler does not play forever; he plays until something happens --- until he is up a hundred dollars,
or broke, or the bar closes. The time he stops is itself random and depends on the game's history. We
need a clean notion of such a time.

\begin{definition}[Stopping time]
\label{ch28:def-stopping}
A random time $T$ taking values in $\set{0,1,2,\dots} \cup \set{\infty}$ is a \emph{stopping time} for
the filtration $(\mathcal{F}_n)$ if, for every $n$, the event $\set{T = n}$ is determined by
$\mathcal{F}_n$.
\end{definition}

The condition says: by time $n$ you always \emph{know} whether you have already stopped. ``Stop the
first time my fortune hits \$100'' is a stopping time --- you can tell at each moment whether it has
happened. ``Stop one step before the peak'' is \emph{not} a stopping time, because recognizing the peak
requires seeing the future. The rule must be enforceable in real time, with no clairvoyance. This is
the same no-peeking constraint as predictability, and it is the whole reason the next theorem holds.

Recall that $\E[M_n] = \E[M_0]$ for any \emph{fixed} time $n$. The burning question is whether this
survives if we replace the fixed $n$ by a \emph{random} stopping time $T$. Is $\E[M_T] = \E[M_0]$?
If yes, then stopping cleverly cannot help. The optional stopping theorem says: yes, provided we rule
out two specific loopholes.

\begin{theorem}[Optional Stopping Theorem, Doob]
\label{ch28:thm-ost}
Let $(M_n)$ be a martingale and $T$ a stopping time. Then $\E[M_T] = \E[M_0]$ provided any one of the
following holds:
\begin{enumerate}
\item $T$ is bounded: there is a constant $N$ with $T \le N$ almost surely; or
\item $T < \infty$ almost surely and the stopped process $(M_{\min(n,T)})$ is uniformly bounded
      (there is a constant $K$ with $\abs{M_{\min(n,T)}} \le K$ for all $n$); or
\item $\E[T] < \infty$ and the increments are bounded, $\abs{M_{n+1} - M_n} \le c$ for some constant
      $c$.
\end{enumerate}
For a supermartingale (unfavorable game) the corresponding conclusion is the inequality
$\E[M_T] \le \E[M_0]$: stopping cannot raise your expected fortune above where you started.
\end{theorem}

\begin{intuitionbox}
The cleanest version is case 1, and you can almost prove it in your head. ``Play until time $T$, then
stop'' is the same as the betting strategy ``bet one unit each round while $n \le T$, and bet zero
afterward.'' The stake $C_n = \ind\set{T \ge n}$ is predictable --- you can tell at the start of round
$n$ whether you have already stopped --- so by Theorem~\ref{ch28:thm-transform} the stopped fortune is
a martingale, and its expectation is the starting value. The conditions in cases 2 and 3 exist
to handle \emph{unbounded} stopping times, where you must ensure the game does not run away to
infinity before $T$ arrives. That is exactly the loophole the doubling system tries to exploit.
Figure~\ref{ch28:fig-stopping-time} illustrates how a bounded stopping time freezes sample paths at absorbing barriers.
\end{intuitionbox}

\begin{figure}[H]
\centering
\begin{tikzpicture}
\begin{axis}[
    width=12.5cm, height=6.2cm,
    xlabel={step $n$}, ylabel={stopped fortune $M_{T \wedge n}$},
    xmin=0, xmax=22, ymin=-5.5, ymax=5.5,
    ytick={-4,-2,0,2,4},
    yticklabels={$-b=-4$, $-2$, $0$, $+2$, $+a=+4$},
    title style={font=\small},
    title={Martingale sample paths under a stopping time: boundary absorption},
    grid=both, grid style={gray!20},
    legend pos=south east, legend style={font=\footnotesize},
]
% Stopping boundaries
\addplot[black, dashed, very thick, domain=0:22] {4};
\addlegendentry{upper boundary $+a = 4$}
\addplot[black, dashed, very thick, domain=0:22] {-4};
\addlegendentry{lower boundary $-b = -4$}

% Path 1: hits +4 at step 8
\addplot[sourceblue, very thick] coordinates {
  (0,0)(1,1)(2,0)(3,1)(4,2)(5,1)(6,2)(7,3)(8,4)
};
\addplot[sourceblue, thick, dotted] coordinates {
  (8,4)(10,4)(12,4)(14,4)(16,4)(18,4)(20,4)(22,4)
};
\addlegendentry{Path 1 (stops at $T_1=8$)}

% Path 2: hits -4 at step 11
\addplot[trapred, very thick] coordinates {
  (0,0)(1,-1)(2,-2)(3,-1)(4,-2)(5,-1)(6,-2)(7,-1)(8,-2)(9,-3)(10,-3)(11,-4)
};
\addplot[trapred, thick, dotted] coordinates {
  (11,-4)(13,-4)(15,-4)(17,-4)(19,-4)(21,-4)(22,-4)
};
\addlegendentry{Path 2 (stops at $T_2=11$)}

% Path 3: hits +4 at step 14
\addplot[targetgreen, very thick] coordinates {
  (0,0)(1,1)(2,2)(3,1)(4,0)(5,-1)(6,0)(7,1)(8,0)(9,1)(10,2)(11,1)(12,2)(13,3)(14,4)
};
\addplot[targetgreen, thick, dotted] coordinates {
  (14,4)(16,4)(18,4)(20,4)(22,4)
};
\addlegendentry{Path 3 (stops at $T_3=14$)}

% Vertical stopping indicators
\draw[gray, dashed] (axis cs:8,-5.5) -- (axis cs:8,4);
\node[font=\scriptsize, gray, anchor=north] at (axis cs:8,-5.5) {$T_1=8$};

\draw[gray, dashed] (axis cs:11,-5.5) -- (axis cs:11,-4);
\node[font=\scriptsize, gray, anchor=north] at (axis cs:11,-5.5) {$T_2=11$};

\draw[gray, dashed] (axis cs:14,-5.5) -- (axis cs:14,4);
\node[font=\scriptsize, gray, anchor=north] at (axis cs:14,-5.5) {$T_3=14$};

\end{axis}
\end{tikzpicture}
\caption{Martingale sample paths frozen by a stopping time $T = \inf\{n : M_n \in \{-b, +a\}\}$. Each path wanders fairly until hitting either $+a$ or $-b$. Once absorbed, the stopped martingale $M_{T \wedge n}$ freezes (dotted lines). By Doob's Optional Stopping Theorem, the fair game guarantees $\E[M_T] = M_0 = 0$, enforcing the exact hitting probabilities $a\,\P(\text{hit } -b) = b\,\P(\text{hit } +a)$.}
\label{ch28:fig-stopping-time}
\end{figure}

\subsection{The doubling system, finally killed}

Return to the opening fantasy. You bet to recover all losses plus one dollar; you stop, at time $T$, at
the first win. Since red appears eventually, $T < \infty$ almost surely, and at time $T$ your profit is
exactly \$1. So $\E[M_T] = 1$, while $M_0 = 0$. Optional stopping appears to be violated --- has the
gambler beaten the fair game?

\begin{paradoxbox}
\textbf{The martingale betting system.} The profit really is \$1 with probability $1$. So where is the
catch? It hides in the \emph{capital required}. If your first win comes on round $k$, your bets were
$1, 2, 4, \dots, 2^{k-1}$ and your peak loss just before winning was $2^{k-1} - 1$ dollars. The chance
of needing exactly $k$ rounds is $2^{-k}$, so the expected size of your largest required bet is
\[
  \sum_{k=1}^{\infty} 2^{-k}\cdot 2^{k-1} = \sum_{k=1}^{\infty}\tfrac12 = \infty.
\]
Your expected maximum exposure is \emph{infinite}. To guarantee the \$1, you must be willing and able
to lose arbitrarily large amounts. No real gambler has infinite capital, and no real casino has no
table limit. The moment you impose any finite bankroll $B$, the strategy becomes a stopping time that
violates none of the theorem's hypotheses, and the conclusion snaps back into place: $\E[M_T] = 0$. You
win a small amount with high probability and suffer a catastrophic, capital-wiping loss with the small
remaining probability, and these balance \emph{exactly} to zero expected profit. The optional stopping
theorem did not fail; the gambler's hidden assumption of infinite wealth did. This is precisely loophole
case 2 of Theorem~\ref{ch28:thm-ost}: with infinite capital the stopped process is \emph{not}
uniformly bounded, so the guarantee evaporates.
\end{paradoxbox}

Figure~\ref{ch28:fig-doubling} shows what this looks like in practice over many simulated nights at
the table.

\begin{figure}[H]
\centering
\begin{tikzpicture}
\begin{axis}[
    width=12cm, height=6.2cm,
    xlabel={night number}, ylabel={cumulative profit (units)},
    title style={font=\small},
    title={Doubling system with a finite bankroll: many small wins, rare ruin},
    grid=both, grid style={gray!20},
    xmin=0, xmax=60, ymin=-130, ymax=40,
    legend pos=south west, legend style={font=\footnotesize},
]
\addplot[thick, sourceblue] coordinates {
(0,0)(5,5)(10,10)(15,15)(20,20)(24,24)(25,-103)(30,-98)(35,-93)(40,-88)(45,-83)(50,-78)(55,-73)(60,-68)
};
\addlegendentry{one gambler's bankroll path}
\addplot[densely dashed, trapred, thick] coordinates {(0,0)(60,0)};
\addlegendentry{break-even}
\end{axis}
\end{tikzpicture}
\caption{A typical realization of the doubling system under a finite bankroll. The fortune climbs
steadily by one unit per night --- the seductive part --- until a long losing streak exceeds the
bankroll and wipes out all accumulated gains and more. Average it over all possible paths and the
expected profit is exactly zero, as the optional stopping theorem demands.}
\label{ch28:fig-doubling}
\end{figure}

\section{The gambler's ruin, in one line}
\label{ch28:gamblers-ruin}

Optional stopping is not merely a no-go theorem; it is also a sharp \emph{computational} tool. The
classic gambler's ruin problem normally requires solving a difference equation. Martingales reduce it
to arithmetic.

\begin{example}[Symmetric gambler's ruin]
\label{ch28:ex-ruin}
You start with \$$a$ and play fair \$1 bets until you either reach a target \$$b$ (with $0 < a < b$)
or go broke at \$0. Let $S_n$ be your fortune and $T$ the first time you hit $0$ or $b$. What is the
probability $p = \P(\text{reach } b \text{ before } 0)$?

The fortune $S_n$ is a martingale (Example~\ref{ch28:ex-walk}, shifted to start at $a$). The stopped
process stays in $[0,b]$, hence is bounded, and $T < \infty$ almost surely (the walk cannot avoid the
boundaries forever). So optional stopping case 2 applies:
\[
  a = \E[S_0] = \E[S_T] = b \cdot p + 0\cdot(1-p) = b\,p
  \quad\Longrightarrow\quad p = \frac{a}{b}.
\]
That is the entire derivation. The probability of reaching your target before ruin is just the ratio of
your starting stake to your target. Want to turn \$100 into \$200 at a fair game? Your chance is exactly
$1/2$. Want to turn \$100 into \$10{,}000? Exactly $1/100$. The deck is not stacked against you --- the
game is perfectly fair --- yet ambitious targets are punishingly unlikely, simply because you have so
much more room to fall to zero than to climb to glory.
\end{example}

Figure~\ref{ch28:fig-ruin} plots this success probability against the size of your target, contrasting
the perfectly fair game with one carrying a tiny house edge. The fair curve decays gently as
$a/b$; the merest bias bends it into a cliff.

\begin{figure}[H]
\centering
\begin{tikzpicture}
\begin{axis}[
    width=12cm, height=6.4cm,
    xlabel={target \$$b$ (starting from \$$a=10$)},
    ylabel={$\P(\text{reach } b \text{ before } 0)$},
    title style={font=\small},
    title={Gambler's ruin: probability of success versus the ambition of the target},
    grid=both, grid style={gray!20},
    xmin=10, xmax=200, ymin=0, ymax=1,
    legend pos=north east, legend style={font=\footnotesize},
    samples=80,
]
% fair game: p = a/b with a = 10
\addplot[thick, sourceblue, domain=10:200] {10/x};
\addlegendentry{fair game, $p=a/b$}
% biased against you: ((q/p)^a - 1)/((q/p)^b - 1) with q/p = 51/49, a = 10
\addplot[thick, trapred, domain=10:200, samples=120] {((51/49)^10 - 1)/((51/49)^x - 1)};
\addlegendentry{$1\%$ edge against you}
\end{axis}
\end{tikzpicture}
\caption{Starting from \$$10$, the chance of reaching a target \$$b$ before going broke. At a perfectly
fair game (blue) the probability is exactly the ratio $a/b$, decaying gently as $1/b$. Introduce a mere
one-percent house edge (red, $q/p = 51/49$) and the success probability collapses exponentially in the
distance to the target --- the same arithmetic that keeps casinos solvent. Both curves come straight
from optional stopping applied to the appropriate martingale, with no difference equations.}
\label{ch28:fig-ruin}
\end{figure}

\begin{example}[How long does ruin take?]
\label{ch28:ex-duration}
Now use the variance martingale of Example~\ref{ch28:ex-variance}. Let $D_n = S_n - a$ be the
displacement from the start, so $(D_n^2 - n)$ is a martingale started at $0$, and optional stopping
gives $\E[D_T^2 - T] = 0$, that is,
\[
  \E[T] = \E[D_T^2].
\]
At time $T$ the displacement from the start is either $-a$ (ruin, probability $1-p = 1 - a/b$) or
$b - a$ (success, probability $p = a/b$). Therefore
\[
  \E[T] = a^2\left(1 - \frac{a}{b}\right) + (b-a)^2\,\frac{a}{b}
        = a^2 \cdot \frac{b-a}{b} + (b-a)^2 \cdot \frac{a}{b}
        = \frac{a(b-a)}{b}\big[a + (b-a)\big] = a(b-a).
\]
So the expected number of bets is $\E[T] = a(b-a)$. With \$$50$ trying to reach \$$100$, you expect
$50\cdot 50 = 2500$ bets before the matter is settled. Fair games are not quick.
\end{example}

\begin{example}[Biased ruin via an exponential martingale]
\label{ch28:ex-biased}
What if the game is unfair, $\P(X_i = +1) = p \ne \tfrac12$, $q = 1-p$? The walk itself is no longer a
martingale, but a cleverly chosen transformation is. Consider $M_n = (q/p)^{S_n}$. Then
\[
  \E[M_{n+1}\given\mathcal{F}_n] = (q/p)^{S_n}\Big[p\cdot(q/p) + q\cdot(p/q)\Big]
  = (q/p)^{S_n}\big[q + p\big] = M_n,
\]
so $(q/p)^{S_n}$ is a martingale. Optional stopping with starting fortune $a$ gives
$(q/p)^a = \E[M_T] = p_{\text{win}}(q/p)^b + (1-p_{\text{win}})(q/p)^0$, and solving,
\[
  \P(\text{reach } b \text{ before } 0) = \frac{(q/p)^a - 1}{(q/p)^b - 1}.
\]
Plug in $p = 18/38$ for an even-money bet on red at American roulette and you will see why even a tiny
house edge, compounded over a long climb to a distant target, makes success vanishingly unlikely. The
martingale method delivers this famous formula with no difference equations at all.
\end{example}

\begin{applicationbox}
\textbf{Finance: fair prices and the absence of arbitrage.} The same logic that forbids beating a
casino is the foundation of modern option pricing. In the risk-neutral framework, a discounted asset
price $\tilde{S}_n$ is modeled as a martingale: $\E[\tilde{S}_{n+1}\given\mathcal{F}_n] =
\tilde{S}_n$. A trading strategy is a predictable sequence of holdings $C_n$, and your discounted
wealth is exactly the martingale transform $\sum C_k(\tilde{S}_k - \tilde{S}_{k-1})$. By
Theorem~\ref{ch28:thm-transform} its expectation stays put --- which is precisely the statement that
\emph{no self-financing strategy can manufacture a guaranteed profit from nothing}: no arbitrage. The
fair price of an option that pays $Z$ at time $N$ is then its discounted risk-neutral expectation
$\E[\tilde{Z}]$, computable by the same optional-stopping machinery. The Black--Scholes formula is the
continuous-time limit of exactly this idea. ``You cannot beat a fair game'' and ``you cannot get free
money in an efficient market'' are the \emph{same theorem}.
\end{applicationbox}

\section{Martingale convergence: where it all settles}
\label{ch28:convergence}

A martingale has constant mean, but it can fluctuate. Does it eventually settle down? The answer is one
of the central theorems in probability, and it asks for almost no hypotheses. A martingale that cannot
grow too positive must converge.

\begin{theorem}[Martingale Convergence Theorem, Doob]
\label{ch28:thm-convergence}
If $(M_n)$ is a martingale (or even a supermartingale) bounded in $L^1$, meaning
$\sup_n \E\abs{M_n} < \infty$, then $M_n$ converges almost surely to a finite limit $M_\infty$. In
particular, any \emph{nonnegative} supermartingale converges almost surely (since for it
$\E\abs{M_n} = \E[M_n] \le \E[M_0] < \infty$ automatically).
\end{theorem}

\begin{intuitionbox}
The proof idea is delightful and is called the \emph{upcrossing} argument. Suppose, toward a
contradiction, that $M_n$ failed to converge --- then it would have to oscillate, crossing some
interval $[\alpha,\beta]$ from below to above infinitely often. But each such ``upcrossing'' is a
buy-low sell-high round trip, and Theorem~\ref{ch28:thm-transform} guarantees you cannot profit from a
predictable strategy on a (super)martingale. So you cannot afford infinitely many upcrossings: each one
would be free money, and there is no free money. Doob made this precise with an inequality bounding the
expected number of upcrossings, and finitely many upcrossings of every interval forces convergence.
Once again, ``you cannot beat a fair game'' is doing the real work.
\end{intuitionbox}

Apply this to our gallery. The likelihood ratio $L_n$ (Example~\ref{ch28:ex-lr}) is nonnegative with
mean $1$, so it converges almost surely --- and when the rival hypothesis is wrong, it converges to
$0$, which is exactly why accumulating data eventually crushes the false model.
Figure~\ref{ch28:fig-lr} shows the log-likelihood-ratio sliding steadily toward $-\infty$ (so
$L_n \to 0$) as data piles up against a false alternative.

\begin{figure}[H]
\centering
\begin{tikzpicture}
\begin{axis}[
    width=12cm, height=5.6cm,
    xlabel={number of observations $n$}, ylabel={$\log L_n$},
    title style={font=\small},
    title={Likelihood ratio against a false model: $\log L_n \to -\infty$, so $L_n \to 0$},
    grid=both, grid style={gray!20},
    xmin=0, xmax=50, ymin=-30, ymax=4,
    legend pos=south west, legend style={font=\footnotesize},
]
% truth N(0,1), alternative N(1,1): per-step log-LR = Y - 1/2, mean -1/2 under truth.
\addplot[thick, sourceblue] coordinates {
(0,0)(5,-1.8)(10,-4.5)(15,-6.2)(20,-9.8)(25,-12.1)(30,-15.0)(35,-17.4)(40,-20.9)(45,-22.6)(50,-25.7)};
\addlegendentry{a simulated $\log L_n$ path}
\addplot[densely dashed, trapred, thick, domain=0:50] {-0.5*x};
\addlegendentry{drift $-n/2$ (the KL rate)}
\end{axis}
\end{tikzpicture}
\caption{The log-likelihood ratio for $\Normal(0,1)$ truth versus an $\Normal(1,1)$ alternative. Each
step adds a term of mean $-\tfrac12$ under the truth (this $\tfrac12$ is exactly the Kullback--Leibler
divergence $\KL(f\,\|\,g)$ between the two models), so $\log L_n$ marches off to $-\infty$ and the
likelihood ratio $L_n$ collapses to $0$. The martingale property guarantees $\E[L_n] = 1$ at every
step, yet almost every path still tends to zero --- a striking reminder that a mean-one martingale can
converge to a limit with mean strictly less than one.}
\label{ch28:fig-lr}
\end{figure}

Polya's urn fraction
(Example~\ref{ch28:ex-polya}) is a bounded martingale, so the red proportion converges to a genuine
random limit. Figure~\ref{ch28:fig-polya} shows several urn paths converging to their individual
destinies.

\begin{figure}[H]
\centering
\begin{tikzpicture}
\begin{axis}[
    width=12cm, height=6.2cm,
    xlabel={number of draws $n$}, ylabel={fraction of red balls $M_n$},
    title={Polya's urn: a martingale that converges to a random limit},
    title style={font=\small},
    grid=both, grid style={gray!20},
    xmin=0, xmax=50, ymin=0, ymax=1,
    legend pos=outer north east, legend style={font=\footnotesize},
]
\addplot[thick, sourceblue] coordinates {
(0,0.5)(1,0.667)(2,0.75)(4,0.7)(6,0.75)(9,0.8)(14,0.81)(20,0.83)(30,0.82)(40,0.83)(50,0.84)};
\addlegendentry{path A}
\addplot[thick, targetgreen] coordinates {
(0,0.5)(1,0.333)(2,0.25)(4,0.3)(6,0.28)(9,0.27)(14,0.25)(20,0.24)(30,0.23)(40,0.23)(50,0.22)};
\addlegendentry{path B}
\addplot[thick, quotientorange] coordinates {
(0,0.5)(1,0.667)(3,0.5)(5,0.52)(8,0.5)(12,0.49)(18,0.5)(25,0.51)(35,0.5)(50,0.5)};
\addlegendentry{path C}
\addplot[thick, structurepurple] coordinates {
(0,0.5)(1,0.333)(2,0.5)(4,0.6)(6,0.625)(10,0.65)(16,0.66)(24,0.65)(34,0.66)(50,0.66)};
\addlegendentry{path D}
\end{axis}
\end{tikzpicture}
\caption{Four simulated runs of Polya's urn. Every run is a martingale with the same starting value
$M_0 = 1/2$, yet each converges to a \emph{different} random limit. Averaged across all runs the mean
stays at $1/2$ forever (constant expectation), but individual paths do not converge to $1/2$ --- the limit is
in fact uniformly distributed on $[0,1]$.}
\label{ch28:fig-polya}
\end{figure}

\begin{aestheticbox}
The convergence theorem is the martingale world's answer to the law of large numbers, but it is far
more general. The law of large numbers needs i.i.d.\ (or at least independent) summands; the martingale
convergence theorem needs only the fair-game property and a mild boundedness condition. Dependence,
feedback, even pathological correlation structures are all permitted. This is why martingales appear
wherever a quantity evolves with memory but no systematic drift --- and why they are the natural
language for stochastic calculus, sequential analysis, and the analysis of randomized algorithms. The
single principle ``you cannot profit from a fair game by stopping or strategizing'' grows
into a complete theory of when random processes must settle down.
\end{aestheticbox}

\section{A worked synthesis: the ballot problem}
\label{ch28:ballot}

To see the machinery dance, here is a classic with a one-line martingale solution.

\begin{example}[Bertrand's ballot problem]
\label{ch28:ex-ballot}
In an election candidate A receives $\alpha$ votes and candidate B receives $\beta$ votes, with
$\alpha > \beta$. Ballots are counted in a uniformly random order. What is the probability that A is
\emph{strictly ahead} throughout the entire count?

Read the count \emph{backwards} from the final tally. Let $M_k$ be the running value
(A's lead)$/$(votes counted) when read from the end; this turns out to be a martingale in the reverse
direction (a ``backward martingale,'' the same projection idea run in reverse). Optional stopping at
the first time A's lead would touch zero, combined with the constant-expectation property, collapses the
whole combinatorial problem to the clean answer
\[
  \P(\text{A always strictly ahead}) = \frac{\alpha - \beta}{\alpha + \beta}.
\]
With $\alpha = 6$, $\beta = 4$, the probability A is never tied or behind is $(6-4)/(6+4) = 1/5$. A
problem that looks like it demands a brutal lattice-path count yields to a single martingale identity.
This is the recurring lesson of the chapter: once you spot the fair game hiding inside a problem, the
answer often falls out of constant expectation and optional stopping with almost no computation.
\end{example}

\begin{connectionbox}
Martingales are the discrete skeleton of the continuous theory you will meet in the next chapter on
Brownian motion. Brownian motion $B_t$ is the continuous-time limit of the symmetric random walk and is
itself a martingale; $B_t^2 - t$ is the continuous analogue of Example~\ref{ch28:ex-variance}; and the
optional stopping theorem becomes Doob's optional sampling theorem, the tool that computes exit
probabilities and hitting times for diffusions. The stochastic integral $\int C_t\,\dd B_t$ is the
continuous martingale transform, and Ito's formula is the rule for how functions of a martingale
behave. Everything in this chapter has a continuous twin, and the discrete versions you have just seen
are the cleanest way to build the intuition before the analysis gets technical. (For the geometric view
of conditional expectation as a projection, and for the integration-as-averaging picture behind these
densities, see the companion calculus and linear algebra readers.)
\end{connectionbox}

\section{Exercises}
\label{ch28:exercises}

\begin{exercise}[Warm-up identities]
\label{ch28:exer-warmup}
Let $(S_n)$ be the symmetric simple random walk with $S_0 = 0$. Verify directly from the definition
that each of the following is a martingale: (a) $S_n$; (b) $S_n^2 - n$; (c) $S_n^3 - 3n S_n$. For part
(c) you will need $\E[X] = 0$, $\E[X^2] = 1$, and $\E[X^3] = 0$ for a fair $X = \pm 1$.
\end{exercise}

\begin{exercise}[Reading the definition]
\label{ch28:exer-def}
Classify each process as a martingale, submartingale, or supermartingale, and justify in one sentence:
(a) a gambler's fortune at a fair coin game; (b) a gambler's fortune at roulette; (c) $\abs{S_n}$ for
the symmetric walk; (d) the running maximum $\max_{k\le n} S_k$. \emph{Hint:} $\abs{x}$ is convex, and
conditional Jensen turns a martingale into a submartingale under any convex function.
\end{exercise}

\begin{exercise}[Predictable strategies]
\label{ch28:exer-predictable}
You play a fair $\pm 1$ game. Your friend proposes: ``bet \$1 normally, but after any loss, bet \$2 on
the next round to make up for it.'' Write this as a predictable stake sequence $C_n$ and use
Theorem~\ref{ch28:thm-transform} to compute the expected fortune after $n$ rounds. Explain in words why
no such rule can produce a positive expected gain.
\end{exercise}

\begin{exercise}[Gambler's ruin numbers]
\label{ch28:exer-ruin}
At a fair game you start with \$$30$ and quit when you reach \$$100$ or go broke. (a) Find the
probability you reach \$$100$. (b) Find the expected number of bets you will make. (c) Repeat (a) for a
target of \$$1000$ and comment on how the success probability scales with the ambition of the target.
\end{exercise}

\begin{exercise}[Biased ruin]
\label{ch28:exer-biased}
A game gives you $+\$1$ with probability $p = 0.49$ and $-\$1$ with probability $q = 0.51$. Using the
exponential martingale of Example~\ref{ch28:ex-biased}, find the probability of turning \$$100$ into
\$$110$ before going broke. Compare to the symmetric-game answer $a/b$ and remark on how much a
$1\%$ edge against you costs.
\end{exercise}

\begin{exercise}[Likelihood ratio]
\label{ch28:exer-lr}
Data $Y_i$ are i.i.d.\ $\Normal(0,1)$ under the truth, but you track the likelihood ratio against the
alternative $\Normal(1,1)$. (a) Write $L_n$ explicitly as an exponential of a sum. (b) Show
$\E[L_n] = 1$ for every $n$ by direct computation, confirming the martingale property. (c) Argue using
Example~\ref{ch28:ex-lr} and the convergence theorem that $L_n \to 0$ almost surely, and explain why
this is exactly what a good sequential test should do.
\end{exercise}

\begin{exercise}[Polya's urn limit]
\label{ch28:exer-polya}
For Polya's urn started with one red and one black ball, let $M_n = R_n/(n+2)$ be the red fraction.
(a) Re-derive that $(M_n)$ is a martingale. (b) Compute $\E[M_n]$ for all $n$. (c) Simulate (by hand or
code) several runs to 100 draws and observe that different runs settle to different limits, consistent
with the claim that the limit is $\Unif(0,1)$.
\end{exercise}

\begin{exercise}[Why stopping cannot help --- conceptual]
\label{ch28:exer-stopping}
Explain carefully why ``stop one step before $S_n$ first reaches its all-time maximum'' is \emph{not} a
valid stopping time, and why ``stop the first time $S_n = 10$'' \emph{is}. Then explain in your own
words why the optional stopping theorem would be false if we allowed the first kind of rule. Connect
this to the no-clairvoyance idea behind predictable strategies.
\end{exercise}

\begin{exercise}[A proof to do yourself]
\label{ch28:exer-proof}
Prove that if $(M_n)$ is a martingale and $\phi$ is a convex function with $\E\abs{\phi(M_n)} < \infty$
for all $n$, then $(\phi(M_n))$ is a submartingale. \emph{Hint:} apply the conditional form of Jensen's
inequality, $\E[\phi(M_{n+1})\given\mathcal{F}_n] \ge \phi(\E[M_{n+1}\given\mathcal{F}_n])$, and then
use the martingale property. Deduce that $S_n^2$ and $\abs{S_n}$ are submartingales for the symmetric
walk.
\end{exercise}

\begin{exercise}[Optional stopping in finance]
\label{ch28:exer-finance}
A discounted stock price is modeled as a martingale with $\tilde{S}_0 = 100$. A trader claims a
strategy that, by buying and selling based only on past prices, guarantees an expected discounted
profit of \$$5$. Using Theorem~\ref{ch28:thm-transform}, explain why the trader must be mistaken (or
secretly using future information). What real-world feature would have to break for a genuine guaranteed
profit --- an arbitrage --- to exist?
\end{exercise}
